\section{Experiments and Results}\label{sec:results}

\subsection{Experimental Setup}
Although standard network simulators such as ns-3, Veins, or SUMO (Simulation of Urban MObility) provide comprehensive full-stack emulation, evaluating federated multi-agent reinforcement learning across heterogeneous RATs requires tight, synchronous coupling between the physical environment and the distributed training loop. We therefore implemented a custom Python-based V2X simulator tailored for algorithmic evaluation and fast DRL inference. To ensure fidelity, the physical-layer behavior rigorously follows established standards, including 3GPP TR 38.901 for cellular fading and IEEE 802.11ax for WiFi. This design choice emphasizes the validation of the multi-agent decision surface and offloading orchestration rather than introducing a new generalized network simulator.



Experiments are conducted in a custom Python V2X simulator. The primary comparison uses N=32 vehicles under a URLLC stress scenario, with five independent random seeds. We retain N=8 results from an earlier controlled evaluation as a reference point and use the fleet density sweep N in \{2, 4, 8, 16, 20, 32\} to evaluate scaling. Vehicle speeds are drawn from $\mathcal{U}[0, 30]$ m/s. Task types follow the URLLC-dominant distribution: 65\% safety (100 ms deadline), 20\% infotainment (500 ms), and 15\% compute (evaluated at a stricter 1000 ms deadline in the experimental results).

Latency measurement scope for \Cref{tab:deadline_analysis}. The per-task-type deadline analysis uses the same N=8, 5-seed, 5-episode protocol and corrected-physics checkpoints as the rest of this evaluation. Per-PDU end-to-end latency is $L_{total}$ computed according to the model in Section III-D; deadline compliance uses $\text{OTR} = n_{rx\_on\_time}/n_{tx}$ per task deadline $D_i$. Mean latency is computed over delivered PDUs only.


\subsection{Performance Metrics and Baseline Comparison}

The performance metrics focus on deadline compliance and reliability compared against standard baselines.

\subsection{Packet-Level Statistics Across V2X Interactions}

\Cref{tab:packet_stats} presents per-vehicle packet statistics for a representative 100-step episode (8 vehicles, URLLC, 40 km/h). Columns: offload PDU transmission attempts ($n_{tx}$), successful deliveries ($n_{rx}$), PDUs received within task deadline ($n_{rx\_on\_time}$), empirical $\text{PRR} = n_{rx}/n_{tx}$ (\Cref{eq:prr}), and on-time rate = $n_{rx\_on\_time}/n_{tx}$. The TOTAL row aggregates across all vehicles.

\begin{table}[t]
\caption{Packet Statistics: Per-Vehicle Output (8 vehicles, 100-step episode, 40 km/h, URLLC)}
\label{tab:packet_stats}
\centering
% booktabs-restyled
\scriptsize
\setlength{\tabcolsep}{2pt}
\renewcommand{\arraystretch}{1.15}
\begin{tabularx}{\columnwidth}{>{\hsize=1.455\hsize\raggedright\arraybackslash}X>{\hsize=0.909\hsize\raggedleft\arraybackslash}X>{\hsize=0.909\hsize\raggedleft\arraybackslash}X>{\hsize=0.909\hsize\raggedleft\arraybackslash}X>{\hsize=0.909\hsize\raggedleft\arraybackslash}X>{\hsize=0.909\hsize\raggedleft\arraybackslash}X}
\toprule
\textbf{Agent} & \textbf{n\_tx} & \textbf{n\_rx} & \textbf{n\_rx\_on\_\allowbreak time} & \textbf{PRR (n\_rx/n\_tx)} & \textbf{On-Time \%} \\
\midrule
Vehicle 0 & 124 & 123 & 118 & 0.9919 & 95.2\% \\
Vehicle 1 & 118 & 117 & 110 & 0.9915 & 93.2\% \\
Vehicle 2 & 131 & 130 & 124 & 0.9924 & 94.7\% \\
Vehicle 3 & 109 & 108 & 101 & 0.9908 & 92.7\% \\
Vehicle 4 & 126 & 125 & 120 & 0.9921 & 95.2\% \\
Vehicle 5 & 115 & 114 & 107 & 0.9913 & 93.0\% \\
Vehicle 6 & 122 & 121 & 116 & 0.9918 & 95.1\% \\
Vehicle 7 & 119 & 118 & 112 & 0.9916 & 94.1\% \\
\midrule
TOTAL     & 964 & 956 & 908 & 0.9917 & 94.2\% \\
\bottomrule
\end{tabularx}
\end{table}


\Cref{tab:packet_stats} (per-vehicle) reports aggregate packet results across all task types. \Cref{tab:deadline_analysis} gives a compact deep safety-validation summary; raw counts are omitted from the table body to keep the IEEE two-column layout readable.

\begin{table}[h!]
\caption{Compact Per-Message-Type Deadline Analysis (N=8, corrected physics, MA-ROPE full policy, 5 seeds)}
\label{tab:deadline_analysis}
\centering
\resizebox{\columnwidth}{!}{%
\begin{tabular}{|l|r|r|r|r|r|}
\hline
\textbf{Task scope} & \textbf{DL} & \textbf{PRR} & \textbf{OTR} & \textbf{Gap} & \textbf{Mean L} \\
\hline
All tasks & - & 0.9415 & 0.8529 & 8.9\% & 214.5 \\
Safety & 100 & 0.9435 & 0.8243 & 11.9\% & 107.4 \\
Infotainment & 500 & 0.9317 & 0.8752 & 5.7\% & 229.8 \\
Compute & 1000 & 0.9431 & 0.9431 & 0.0\% & 665.6 \\
\hline
\end{tabular}%
}
\end{table}


\Cref{tab:deadline_analysis} should not be compared directly with the N=32 multi-method safety table: it is a MA-ROPE-only, N=8 evaluation used to isolate deadline misses by task type.

The per-task-type breakdown illustrates the PRR-vs-OTR distinction clearly. Safety PDUs (100 ms deadline) achieve PRR = 0.9435 but OTR = 0.8243, leaving a safety-only PRR-OTR gap of 11.9 percentage points. Infotainment PDUs (500 ms deadline) show a smaller gap (5.7 percentage points), and compute PDUs (1000 ms deadline) show none, since their deadline is generous relative to typical local and offload completion times. Fleet-wide, all tasks together achieve PRR = 0.9415 and OTR = 0.8529, an 8.9 percentage point gap. This all-tasks figure is a different quantity from the N=32 fleet-wide values reported elsewhere in this paper and should not be compared to them directly.
